\chapter{Formes hermitiennes}\label{ch:b2:hermitian}

Les espaces vectoriels complexes possèdent leur propre géométrie de
produit scalaire, avec une torsion~: linéarité en une variable,
semi-linéarité en l'autre. Le prix à payer pour accepter cette torsion
est une théorie spectrale encore plus limpide que la théorie réelle
--- les \omterm{def:b2:hermitian:adjoint}{endomorphismes hermitiens} ont des \omterm{def:b2:reduction:eigen}{valeurs propres} réelles, les
endomorphismes \omterm{def:b2:hermitian:adjoint}{unitaires} ont des \omterm{def:b2:reduction:eigen}{valeurs propres} de module~$1$, et les
deux se diagonalisent dans des bases orthonormées. Ce court chapitre
applique le programme euclidien du \cref{ch:b2:quadratic} sur $\C$.

\section{Produits scalaires hermitiens}

\begin{definition}\label{def:b2:hermitian:def}
Un \emph{produit scalaire hermitien}\index{produit scalaire hermitien}
sur un espace vectoriel complexe $E$ est une application
$\langle\cdot,\cdot\rangle \colon E \times E \to \C$ linéaire en la
seconde variable, à symétrie hermitienne ($\langle y, x\rangle =
\conj{\langle x, y\rangle}$ --- donc semi-linéaire en la première
variable), et définie positive ($\langle x, x\rangle > 0$ pour $x \neq
0$). L'exemple standard sur $\C^n$~:
\[
\langle x, y \rangle = \sum_{i=1}^{n} \conj{x_i}\, y_i ;
\]
sur les fonctions \omterm{def:b2:metric:continuity}{continues}, $\langle f, g\rangle = \int_a^b \conj
f\,g$. \omterm{def:b2:nvs:norm}{Norme}~: $\norm x = \sqrt{\langle x,x\rangle}$~; un espace
complexe de dimension finie ainsi muni est un \emph{espace
hermitien}\index{espace hermitien}.
\end{definition}

\begin{example}[Premiers calculs]\label{ex:b2:hermitian:firstcomp}
Dans $\C^2$, prenons $x = (1+\iu,\ 2-\iu)$ et $y = (\iu,\ 1)$.
Alors
\[
\norm x^2 = \abs{1+\iu}^2 + \abs{2-\iu}^2 = 2 + 5 = 7,
\qquad
\norm y^2 = 1 + 1 = 2,
\]
et, en conjuguant le premier argument,
\[
\langle x, y\rangle
= \conj{(1+\iu)}\,\iu + \conj{(2-\iu)}\cdot1
= (1-\iu)\iu + (2+\iu)
= (\iu + 1) + (2 + \iu) = 3 + 2\iu .
\]
Cauchy--Schwarz est vérifiée~: $\abs{\langle x, y\rangle}^2 = 9 +
4 = 13 \leq 14 = \norm x^2\norm y^2$ --- proche de l'égalité,
car $x$ est proche d'un multiple de $y$. Remarquons aussi
$\langle y, x\rangle = \conj{3 + 2\iu} = 3 - 2\iu$~: la symétrie
hermitienne en action, et la raison pour laquelle $\langle x,
x\rangle$ est toujours réel.
\end{example}

\begin{theorem}[Cauchy--Schwarz, cas complexe]\label{thm:b2:hermitian:cs}
$\abs{\langle x, y\rangle} \leq \norm x \norm y$, avec égalité si et
seulement si $x, y$ sont linéairement dépendants~; $\norm\cdot$ est une
\omterm{def:b2:nvs:norm}{norme}. De plus, il existe des bases orthonormées (Gram--Schmidt
s'applique mot pour mot), avec
\[
x = \sum_i \langle e_i, x\rangle\, e_i,
\qquad
\norm x^2 = \sum_i \abs{\langle e_i, x\rangle}^2 .
\]
\end{theorem}

\begin{proof}
Pour $y \neq 0$ et $t \in \C$~: $0 \leq \norm{x - ty}^2 = \norm x^2
- 2\Re\bigl(\conj t\langle y, x\rangle\bigr) + \abs t^2\norm y^2$.
Choisissons $t = \frac{\langle y, x\rangle}{\norm y^2}$~:
\[
0 \leq \norm x^2 - \frac{\abs{\langle y, x\rangle}^2}{\norm y^2},
\]
ce qui est l'inégalité~; l'égalité impose $x = ty$. Inégalité
triangulaire, en entier~:
\[
\norm{x + y}^2 = \norm x^2 + 2\,\Re\langle x, y\rangle
+ \norm y^2
\leq \norm x^2 + 2\,\abs{\langle x, y\rangle} + \norm y^2
\leq \bigl(\norm x + \norm y\bigr)^2 ,
\]
en utilisant $\Re z \leq \abs z$ puis Cauchy--Schwarz~;
l'homogénéité et la séparation sont immédiates, donc $\norm\cdot$ est
une \omterm{def:b2:nvs:norm}{norme}. Gram--Schmidt~: comme dans le cas réel, avec les
conjugaisons placées selon la définition (attention à la
convention~: nos produits sont semi-linéaires en la \emph{première}
variable, donc les coordonnées sont $\langle e_i, x\rangle$~;
l'exemple suivant déroule l'algorithme une fois en entier).
\end{proof}

\begin{example}[Gram--Schmidt complexe, déroulé en entier]\label{ex:b2:hermitian:gramschmidt}
Orthonormalisons la base $v_1 = (1, \iu)$, $v_2 = (0, 1)$ de
$\C^2$. Premier vecteur~: $\norm{v_1}^2 = \abs1^2 + \abs\iu^2 = 2$,
donc $e_1 = \frac{1}{\sqrt2}(1, \iu)$. Projetons $v_2$ --- avec le
conjugué dans la première variable~:
\[
\langle e_1, v_2\rangle
= \frac{1}{\sqrt2}\bigl(\conj{1}\cdot0 +
\conj{\iu}\cdot1\bigr)
= \frac{-\iu}{\sqrt2} ,
\qquad
v_2 - \langle e_1, v_2\rangle e_1
= (0,1) + \frac{\iu}{2}\,(1, \iu)
= \Bigl(\frac\iu2,\ \frac12\Bigr) .
\]
Sa \omterm{def:b2:nvs:norm}{norme} vaut $\sqrt{\frac14 + \frac14} = \frac{1}{\sqrt2}$~:
$e_2 = \frac{1}{\sqrt2}(\iu, 1)$. Vérification~: $\langle e_1,
e_2\rangle = \frac12(\conj1\cdot\iu + \conj\iu\cdot1) =
\frac12(\iu - \iu) = 0$. Coordonnées de $v_2$ dans la nouvelle base~:
$v_2 = \langle e_1, v_2\rangle e_1 + \langle e_2, v_2\rangle
e_2$ avec $\langle e_2, v_2\rangle = \frac{1}{\sqrt2}$ ---
attention à l'\omterm{def:b2:structures:generated}{ordre}~: $\langle v_2, e_1\rangle$ donnerait le
coefficient \emph{conjugué}. Idée finale~: l'algorithme est
celui du cas euclidien mot pour mot~; le seul piège est l'endroit
où tombe la conjugaison, et le calcul de $\norm{v_2 -
\text{proj}}^2 > 0$ utilise silencieusement la positivité --- l'axiome
qui fait fonctionner toute la géométrie.
\end{example}

\section{Adjoint, endomorphismes hermitiens et unitaires}

\begin{definition}\label{def:b2:hermitian:adjoint}
L'\emph{adjoint} $u^*$ de $u \in \mathcal{L}(E)$ est défini par
$\langle u^*(x), y\rangle = \langle x, u(y)\rangle$~; dans une base
orthonormée, $\operatorname{Mat}(u^*) = \conj{A}^{\mathsf T} =:
A^{\dagger}$ (transconjuguée) --- en effet, si $B =
(b_{ij})$ est la matrice de $u^*$ dans la base orthonormée
$(e_i)$, alors $b_{ij} = \langle e_i, u^*(e_j)\rangle$, et
l'identité de définition donne
\[
\conj{b_{ij}}
= \langle u^*(e_j), e_i\rangle
= \langle e_j, u(e_i)\rangle = a_{ji} ,
\qquad\text{donc}\qquad
B = \conj{A}^{\mathsf T} .
\] $u$ est
\emph{hermitien}\index{endomorphisme hermitien} lorsque $u^* = u$
($A^\dagger = A$), \emph{unitaire}\index{groupe unitaire} lorsque $u^*u =
\mathrm{id}$ ($A^\dagger A = I$~: le groupe $U(n)$), \emph{normal}
lorsque $u^*u = uu^*$.
\end{definition}

\begin{proposition}\label{prop:b2:hermitian:eigenvalues}
Les \omterm{def:b2:reduction:eigen}{valeurs propres} d'un \omterm{def:b2:hermitian:adjoint}{endomorphisme hermitien} sont \emph{réelles}~;
les \omterm{def:b2:reduction:eigen}{valeurs propres} d'un endomorphisme \omterm{def:b2:hermitian:adjoint}{unitaire} sont de \emph{module
$1$}~; dans les deux cas, les sous-espaces propres associés à des
\omterm{def:b2:reduction:eigen}{valeurs propres} distinctes sont orthogonaux.
\end{proposition}

\begin{proof}
Cas \omterm{def:b2:hermitian:adjoint}{hermitien}, $u(x) = \lambda x$, $x \neq 0$~:
\[
\lambda \norm x^2 = \langle x, u(x)\rangle = \langle u(x), x\rangle
= \conj{\langle x, u(x)\rangle} = \conj\lambda\,\norm x^2 ,
\]
donc $\lambda \in \R$. Cas \omterm{def:b2:hermitian:adjoint}{unitaire}~: $\norm{u(x)} = \norm x$ (à
partir de $u^*u = \mathrm{id}$), donc $\abs\lambda\norm x = \norm x$.
Orthogonalité (cas \omterm{def:b2:hermitian:adjoint}{hermitien})~: $\lambda\langle x, y\rangle = \langle
u(x), y\rangle = \langle x, u(y)\rangle = \mu\langle x, y\rangle$ pour
des vecteurs propres avec $\lambda \neq \mu$ réels. Cas \omterm{def:b2:hermitian:adjoint}{unitaire}, en
entier~: pour $u(x) = \lambda x$, $u(y) = \mu y$ avec $\lambda \neq
\mu$ (tous deux de module~$1$),
\[
\langle x, y\rangle = \langle u(x), u(y)\rangle
= \conj\lambda\mu\,\langle x, y\rangle ,
\]
et $\conj\lambda\mu = \frac{\mu}{\lambda} \neq 1$~: le facteur
n'est pas $1$, donc $\langle x, y\rangle = 0$.
\end{proof}

\begin{example}[Une matrice anti-hermitienne, diagonalisée]\label{ex:b2:hermitian:skewexample}
$A = \begin{pmatrix} 0 & -2\\ 2 & 0\end{pmatrix}$ vérifie
$A^\dagger = A^{\mathsf T} = -A$~: anti-hermitienne (et aussi réelle
antisymétrique --- sur $\R$ elle n'a aucune \omterm{def:b2:reduction:eigen}{valeur propre}).
\omterm{def:b2:reduction:charpoly}{Polynôme caractéristique} $X^2 + 4$~: \omterm{def:b2:reduction:eigen}{valeurs propres} $\pm2\iu$,
imaginaires pures, comme le prévoit en général
l'\cref{exo:b2:hermitian:9}. Vecteurs propres~: $(A - 2\iu I)v = 0$
donne $v_1 = \frac{1}{\sqrt2}(1, \iu)$, et $v_2 = \frac{1}{\sqrt2}(1,
-\iu)$ pour $-2\iu$~; ils sont orthogonaux~:
\[
\langle v_1, v_2\rangle
= \tfrac12\bigl(\conj{1}\cdot1 +
\conj{\iu}\cdot(-\iu)\bigr)
= \tfrac12(1 - 1) = 0 .
\]
Donc $A = U\operatorname{diag}(2\iu, -2\iu)\,U^\dagger$ avec $U =
(v_1\ v_2)$ \omterm{def:b2:hermitian:adjoint}{unitaire}. Idée finale~: $H = -\iu A =
\begin{pmatrix} 0 & 2\iu\\ -2\iu & 0\end{pmatrix}$ est hermitienne,
de \omterm{def:b2:reduction:eigen}{spectre} réel $\{\pm2\}$ et de \emph{mêmes} vecteurs propres --- la
bijection $u \mapsto \iu u$ entre \omterm{def:b2:hermitian:adjoint}{endomorphismes hermitiens} et
\omterm{def:b2:hermitian:adjoint}{anti-hermitiens} (\cref{exo:b2:hermitian:9}), vue matrice par
matrice~; sur $\R$, la même $A$ est une rotation-homothétie sans
aucun vecteur propre, et seul le passage à $\C$ révèle sa forme
normale.
\end{example}

\begin{theorem}[Théorème spectral hermitien]\label{thm:b2:hermitian:spectral}
Tout \omterm{def:b2:hermitian:adjoint}{endomorphisme hermitien} d'un \omterm{def:b2:hermitian:def}{espace hermitien} admet une base
orthonormée de vecteurs propres (avec des \omterm{def:b2:reduction:eigen}{valeurs propres} réelles)~:
$A^\dagger = A$ implique $A = U D U^{\dagger}$ avec $U \in U(n)$ et $D$
diagonale réelle.\index{théorème spectral!hermitien}
\end{theorem}

\begin{proof}
Sur $\C$, le \omterm{def:b2:reduction:charpoly}{polynôme caractéristique} est scindé~: il existe un
vecteur propre $e_1$ (\cref{ch:b2:reduction}) --- aucun argument de
compacité n'est nécessaire, un avantage de $\C$. Normalisons-le. Son
orthogonal $F = e_1^\perp$ est stable~: pour $x \perp e_1$,
\[
\langle e_1, u(x)\rangle = \langle u(e_1), x\rangle
= \lambda_1\langle e_1, x\rangle = 0
\]
($\lambda_1$ réel). La restriction est hermitienne~; on raisonne par
récurrence sur la dimension et on concatène. En détail~: la
restriction $u|_F$ est un endomorphisme de l'\omterm{def:b2:hermitian:def}{espace hermitien} $F$
(dimension $n - 1$) avec $\langle u|_F(x), y\rangle = \langle x,
u|_F(y)\rangle$ hérité de $u$~; l'hypothèse de récurrence fournit une
base orthonormée $(e_2, \dots, e_n)$ de $F$ formée de vecteurs
propres, et $(e_1, e_2, \dots, e_n)$ est orthonormée dans $E$ ($e_1
\perp F$) et formée de vecteurs propres de $u$. Traduction
matricielle~: les colonnes de $U$ sont les $e_i$, $U^\dagger U = I$
exprime leur orthonormalité, et $AU = UD$ rassemble les équations aux
\omterm{def:b2:reduction:eigen}{valeurs propres}, d'où $A = UDU^\dagger$ avec $D$ diagonale réelle
(\cref{prop:b2:hermitian:eigenvalues}).
\end{proof}

\begin{example}\label{ex:b2:hermitian:pauli}
$A = \begin{pmatrix} 0 & -\iu\\ \iu & 0\end{pmatrix}$ est hermitienne
($A^\dagger = A$)~: \omterm{def:b2:reduction:eigen}{valeurs propres} tirées de $\chi_A = X^2 - 1$~: $\pm
1$ (réelles, comme promis), avec vecteurs propres orthonormés
$\frac{1}{\sqrt2}(1, \iu)^{\mathsf T}$ et $\frac{1}{\sqrt2}(1,
-\iu)^{\mathsf T}$. (Les physiciens connaissent $A$ comme une matrice
de Pauli~; c'est parce que les \omterm{def:b2:reduction:eigen}{spectres} \omterm{def:b2:hermitian:adjoint}{hermitiens} sont réels que les
observables quantiques sont modélisées par des opérateurs \omterm{def:b2:hermitian:adjoint}{hermitiens}.)
\end{example}

\begin{example}[Une matrice hermitienne définie positive, traitée]\label{ex:b2:hermitian:pdrun}
$A = \begin{pmatrix} 2 & 1-\iu\\ 1+\iu & 3\end{pmatrix}$~:
hermitienne, puisque la diagonale est réelle et les coefficients hors
diagonale sont conjugués. \omterm{def:b2:reduction:charpoly}{Polynôme caractéristique}~:
\[
(2-\lambda)(3-\lambda) - \abs{1-\iu}^2
= \lambda^2 - 5\lambda + 4
= (\lambda - 1)(\lambda - 4) :
\]
\omterm{def:b2:reduction:eigen}{spectre} $\{1, 4\}$, réel et positif --- $A$ est définie
positive. Vecteurs propres~: pour $\lambda = 4$, le système $(A -
4I)v = 0$ donne $v_4 = (1 - \iu,\ 2)$ (vérifier la seconde ligne~:
$(1+\iu)(1-\iu) - 2 = 0$)~; pour $\lambda = 1$, $v_1 = (1 - \iu,\
-1)$. Orthogonalité, avec le conjugué dans la première variable~:
\[
\langle v_4, v_1\rangle
= \conj{(1-\iu)}\,(1-\iu) + \conj{2}\,(-1)
= 2 - 2 = 0 . \checkmark
\]
En normalisant ($\norm{v_4}^2 = 2 + 4 = 6$, $\norm{v_1}^2 = 2 + 1
= 3$), on obtient la matrice \omterm{def:b2:hermitian:adjoint}{unitaire} $U = \bigl(\frac{v_4}{\sqrt6}\
\frac{v_1}{\sqrt3}\bigr)$ avec $A =
U\operatorname{diag}(4,1)U^\dagger$. Idée finale~: la lecture de
Rayleigh est immédiate --- sur la sphère unité de
$\C^2$, $\langle x, Ax\rangle$ parcourt $\intcc{1}{4}$,
atteint aux deux vecteurs propres~; c'est le germe en $n = 2$ de la
théorie de Courant--Fischer construite dans le problème du week-end.
Vérifions au passage que la forme est réelle aussi en dehors des
vecteurs propres~: en $x = (1, \iu)$,
\[
Ax = \bigl(2 + (1-\iu)\iu,\ (1+\iu) + 3\iu\bigr)
= (3 + \iu,\ 1 + 4\iu),
\]
\[
\langle x, Ax\rangle = \conj{1}\,(3+\iu) +
\conj{\iu}\,(1+4\iu)
= (3 + \iu) + (-\iu)(1 + 4\iu) = 3 + \iu - \iu + 4 = 7 \in
\R ,
\]
comme le garantit pour tout $x$ le mécanisme de preuve de la
\cref{prop:b2:hermitian:eigenvalues} (symétrie hermitienne contre
$A^\dagger = A$).
\end{example}

\begin{example}[Une matrice unitaire diagonalisée]\label{ex:b2:hermitian:unitary}
$U = \frac{1}{\sqrt2}\begin{pmatrix} 1 & \iu\\ \iu & 1
\end{pmatrix}$ (\omterm{def:b2:hermitian:adjoint}{unitaire} d'après l'\cref{exo:b2:hermitian:2}). Son
\omterm{def:b2:reduction:charpoly}{polynôme caractéristique} est $\bigl(X - \frac{1}{\sqrt2}\bigr)^2
+ \frac12$, de racines
\[
\lambda_\pm = \frac{1 \pm \iu}{\sqrt2} = \eu^{\pm\iu\pi/4},
\]
de module~$1$ comme le promettait la
\cref{prop:b2:hermitian:eigenvalues}, et de vecteurs propres
orthonormés $\frac{1}{\sqrt2}(1, \pm1)$. Donc $U =
V\operatorname{diag}(\eu^{\iu\pi/4},
\eu^{-\iu\pi/4})V^\dagger$~: dans la bonne base, $U$ est une paire de
rotations planes d'angle $\pm\frac\pi4$ --- une matrice de rotation
réelle n'a aucun vecteur propre réel, mais sur $\C$ elle se scinde en
deux scalaires de module~$1$. Idée finale~: les \omterm{def:b2:reduction:eigen}{spectres} \omterm{def:b2:hermitian:adjoint}{hermitiens}
vivent sur la droite réelle, les \omterm{def:b2:reduction:eigen}{spectres} \omterm{def:b2:hermitian:adjoint}{unitaires} sur le cercle
unité~; les deux sont des ombres de la même normalité, et la
transformation de Cayley de l'\cref{exo:b2:hermitian:6} fait passer
d'une image à l'autre.
\end{example}

\begin{example}[La transformation de Cayley, calculée]\label{ex:b2:hermitian:cayleyrun}
Appliquons l'\cref{exo:b2:hermitian:6} à $H = \begin{pmatrix} 0 & 1\\
1 & 0\end{pmatrix}$ (hermitienne, de \omterm{def:b2:reduction:eigen}{spectre} $\{1, -1\}$,
de vecteurs propres orthonormés $\frac{1}{\sqrt2}(1, \pm1)$). Dans la
base propre, tout est scalaire~: la transformation $\lambda
\mapsto \frac{\lambda - \iu}{\lambda + \iu}$ envoie
\[
1 \longmapsto \frac{1 - \iu}{1 + \iu} = -\iu,
\qquad
-1 \longmapsto \frac{-1 - \iu}{-1 + \iu} = \iu ,
\]
(multiplier par le conjugué du dénominateur), donc $U = (H -
\iu I)(H + \iu I)^{-1}$ est la matrice \omterm{def:b2:hermitian:adjoint}{unitaire} de \omterm{def:b2:reduction:eigen}{valeurs propres}
$\mp\iu$ sur ces mêmes vecteurs propres~:
\[
U = \frac{1}{2}\begin{pmatrix} 1 & 1\\ 1 & -1\end{pmatrix}
\begin{pmatrix} -\iu & 0\\ 0 & \iu\end{pmatrix}
\begin{pmatrix} 1 & 1\\ 1 & -1\end{pmatrix}
= \begin{pmatrix} 0 & -\iu\\ -\iu & 0\end{pmatrix} .
\]
Vérification~: $U^\dagger U = I$, et $1 \notin \operatorname{Sp}U =
\{\pm\iu\}$, comme le promet la théorie. Idée finale~: la droite
réelle s'envoie sur le cercle unité privé du point $1$ ---
\omterm{def:b2:reduction:eigen}{valeur propre} par \omterm{def:b2:reduction:eigen}{valeur propre}, la transformation de Cayley est
l'homographie de Möbius $\frac{\lambda-\iu}{\lambda+\iu}$, et les
matrices ne font que suivre leurs \omterm{def:b2:reduction:eigen}{spectres}.
\end{example}

\begin{remark}[Pièges classiques]
\emph{(i) Où tombe la barre~:} ce livre conjugue la
\emph{première} variable, donc les coordonnées sont $\langle e_i,
x\rangle$ et $\langle\lambda x, y\rangle = \conj\lambda\langle x,
y\rangle$~; beaucoup de textes conjuguent plutôt la seconde variable
--- traduire avant de comparer les formules, sinon les signes des
$\iu$ deviennent faux silencieusement. \emph{(ii) La \omterm{def:b2:quadratic:def}{polarisation}
complexe est plus forte~:} sur $\C$, si $\langle x, u(x)\rangle = 0$
pour \emph{tout} $x$, alors $u = 0$ (développer $x + y$ et $x + \iu
y$~: les parties réelle et imaginaire de $\langle x, u(y)\rangle$
s'annulent toutes deux)~; sur $\R$, cela échoue --- la rotation
d'angle $\frac\pi2$ vérifie $\langle x, u(x)\rangle = 0$ partout. Par
conséquent, sur $\C$ uniquement, «~$\langle x, u(x)\rangle \in \R$
pour tout $x$~» impose déjà que $u$ soit \omterm{def:b2:hermitian:adjoint}{hermitien}. \emph{(iii) Une
matrice réelle normale n'est pas \omterm{def:b2:reduction:diag}{diagonalisable}~:} la matrice de
l'\cref{ex:b2:hermitian:skewexample} est normale mais n'a aucune
\omterm{def:b2:reduction:eigen}{valeur propre} réelle~; la diagonalisation \omterm{def:b2:hermitian:adjoint}{unitaire} est un théorème
\emph{sur} $\C$, et sur $\R$ on n'obtient que des réductions par
blocs. \emph{(iv) Vérifier l'unitarité~:} $U^\dagger U = I$ signifie
que les \emph{colonnes} sont orthonormées pour le produit \omterm{def:b2:hermitian:adjoint}{hermitien}
--- tester $UU^{\mathsf T}$, ou oublier la conjugaison dans les
produits de colonnes, sont les deux façons classiques de certifier
une mauvaise matrice.
\end{remark}

\begin{example}[Les isométries sont exactement les applications unitaires]\label{ex:b2:hermitian:isometry}
La conservation de la \omterm{def:b2:nvs:norm}{norme} paraît plus faible que l'unitarité, mais
sur $\C$ elle ne l'est pas~: si $\norm{u(x)} = \norm x$ pour tout
$x$, alors $u^*u = \mathrm{id}$. En effet, $v = u^*u - \mathrm{id}$
est \omterm{def:b2:hermitian:adjoint}{hermitien} et vérifie $\langle x, v(x)\rangle = \norm{u(x)}^2 -
\norm x^2 = 0$ pour tout $x$~; par \omterm{def:b2:quadratic:def}{polarisation} complexe (piège (ii)
ci-dessus), une application dont la «~diagonale~» est identiquement
nulle est nulle~: $v = 0$. Concrètement, la \omterm{def:b2:quadratic:def}{polarisation} s'écrit
\[
0 = \langle x + y, v(x+y)\rangle
= \langle x, v(y)\rangle + \langle y, v(x)\rangle,
\qquad
0 = \langle x + \iu y, v(x + \iu y)\rangle
= \iu\langle x, v(y)\rangle - \iu\langle y, v(x)\rangle ,
\]
et les deux lignes ensemble imposent $\langle x, v(y)\rangle = 0$
pour tous $x, y$. Idée finale~: c'est pourquoi «~\omterm{def:b2:hermitian:adjoint}{unitaire}~» peut se
vérifier en mesurant seulement des longueurs --- une rigidité que le
chapitre de Fourier exploitera, où conserver l'énergie
$\norm f_2$ (Parseval) revient à conserver tous les produits
scalaires des coefficients.
\end{example}

\begin{remark}[Perspectives dans ce volume]
La machinerie hermitienne construite ici est consommée presque
immédiatement. Le chapitre de Fourier \emph{est} de la géométrie
hermitienne en dimension infinie~: les exponentielles $(e_n)$ forment
une famille orthonormée pour $\langle f, g\rangle =
\frac{1}{2\pi}\int\conj fg$, l'inégalité de Bessel est l'estimation de
projection de la \cref{thm:b2:hermitian:cs} de ce chapitre, et
Parseval en est l'égalité limite. La transformation de Fourier finie
(\cref{exo:b2:hermitian:10}) réapparaît chaque fois qu'il faut
diagonaliser une convolution. Et le problème du week-end de ce
chapitre --- Courant--Fischer, Weyl, entrelacement --- fournit la
stabilité des \omterm{def:b2:reduction:eigen}{valeurs propres} qu'invoque le chapitre d'équations
différentielles lorsqu'il affirme que de petites perturbations d'un
système ne déplacent que légèrement ses fréquences. En arrière, tout
ici est le miroir complexe du chapitre sur les \omterm{def:b2:quadratic:def}{formes quadratiques}~:
garder les deux dictionnaires côte à côte ($A^{\mathsf T}
\leftrightarrow A^\dagger$, orthogonal $\leftrightarrow$ \omterm{def:b2:hermitian:adjoint}{unitaire},
Rayleigh réel dans les deux).
\end{remark}

\begin{remark}[Endomorphismes normaux]
Sur $\C$, l'énoncé définitif est~: $u$ est \emph{unitairement
\omterm{def:b2:reduction:diag}{diagonalisable} si et seulement s'il est normal} ($u^*u = uu^*$) ---
ce qui couvre d'un coup les applications hermitiennes, \omterm{def:b2:hermitian:adjoint}{unitaires} et
anti-hermitiennes. La preuve est un agréable renforcement de
l'argument ci-dessus (\cref{exo:b2:hermitian:8}). Sur $\R$, en
revanche, la normalité ne procure qu'une diagonalisation par blocs
(blocs de rotation)~: la géométrie complexe est réellement plus
simple.
\end{remark}

\begin{remark}[Où cela sert]
La théorie spectrale hermitienne est la mathématique de la mécanique
quantique~: les observables sont modélisées par des opérateurs
\omterm{def:b2:hermitian:adjoint}{hermitiens} (\omterm{def:b2:reduction:eigen}{spectres} réels = valeurs mesurables), l'évolution
temporelle par des opérateurs \omterm{def:b2:hermitian:adjoint}{unitaires} (conservation de la \omterm{def:b2:nvs:norm}{norme} =
conservation de la probabilité). Dans ce livre, le chapitre de
Fourier repose sur l'orthonormalité des exponentielles --- un énoncé
de \omterm{def:b2:hermitian:def}{produit scalaire hermitien} --- et la diagonalisation des matrices
circulantes (\cref{exo:b2:hermitian:10}) est la transformation de
Fourier finie. Le problème du week-end développe le calcul
variationnel des \omterm{def:b2:reduction:eigen}{valeurs propres} (Courant--Fischer, Weyl,
entrelacement), le pain quotidien de l'analyse numérique et de la
physique mathématique~; le volume de l'année~3 l'étend aux opérateurs
autoadjoints \omterm{def:b2:metric:compact}{compacts} sur les espaces de Hilbert.
\end{remark}

\section{Exercices}

\begin{exercise}[$\star$]\label{exo:b2:hermitian:1}
Sur $\C^2$~: calculer $\langle x, y\rangle$, $\norm x$, $\norm y$ pour
$x = (1, \iu)$, $y = (\iu, 1)$~; sont-ils orthogonaux~? Donner une
base orthonormée contenant $\frac{x}{\norm x}$.
\end{exercise}

\begin{exercise}[$\star$]\label{exo:b2:hermitian:2}
Lesquelles sont hermitiennes~? \omterm{def:b2:hermitian:adjoint}{unitaires}~? normales~?
\[
\begin{pmatrix} 1 & \iu\\ -\iu & 2 \end{pmatrix},
\qquad
\frac{1}{\sqrt2}\begin{pmatrix} 1 & \iu\\ \iu & 1\end{pmatrix},
\qquad
\begin{pmatrix} 0 & 1\\ 0 & 0 \end{pmatrix}.
\]
\end{exercise}

\begin{exercise}[$\star$]\label{exo:b2:hermitian:3}
Montrer qu'une matrice $A \in \mathcal{M}_n(\C)$ s'écrit de manière
unique $A = H + \iu K$ avec $H, K$ hermitiennes \emph{(les «~parties
réelle et imaginaire~» $H = \frac{A + A^\dagger}{2}$, $K = \frac{A -
A^\dagger}{2\iu}$)}, et que $A$ est normale si et seulement si $H$ et
$K$ commutent.
\end{exercise}

\begin{exercise}[$\star\star$]\label{exo:b2:hermitian:4}
Diagonaliser dans une base orthonormée~: $A = \begin{pmatrix} 2 & \iu\\
-\iu & 2\end{pmatrix}$, et calculer $A^k$ pour $k \in \N$.
\end{exercise}

\begin{exercise}[$\star\star$]\label{exo:b2:hermitian:5}
Montrer que $U(n)$ est \omterm{def:b2:metric:compact}{compact}, et que l'application \omterm{def:b2:reduction:eigen}{valeur propre} est
surjective~: tout $\lambda$ de module~$1$ apparaît pour une certaine
matrice \omterm{def:b2:hermitian:adjoint}{unitaire}. Montrer que $\det U \in \mathbb{U}$ (le cercle
unité) pour $U \in U(n)$.
\end{exercise}

\begin{exercise}[$\star\star$]\label{exo:b2:hermitian:6}
(Transformation de Cayley) Soit $H$ hermitienne. Montrer que $H + \iu
I$ est inversible et que $U = (H - \iu I)(H + \iu I)^{-1}$ est
\omterm{def:b2:hermitian:adjoint}{unitaire}, avec $1 \notin \operatorname{Sp}(U)$. \emph{(Travailler
spectralement~: sur une base propre de $H$, tout est scalaire.)}
\end{exercise}

\begin{exercise}[$\star\star$]\label{exo:b2:hermitian:7}
Pour $A$ hermitienne définie positive ($\langle x, Ax\rangle > 0$ pour
$x \neq 0$), montrer que $\operatorname{Sp}(A) \subseteq
\intoo{0}{\infty}$, que $A = B^2$ pour une certaine matrice $B$
hermitienne définie positive, et que $\det A > 0$.
\end{exercise}

\begin{exercise}[$\star\star\star$]\label{exo:b2:hermitian:8}
(Théorème spectral pour les endomorphismes normaux) Soit $u$ normal
sur un \omterm{def:b2:hermitian:def}{espace hermitien}.
\begin{enumerate}
  \item Montrer $\norm{u(x)} = \norm{u^*(x)}$ pour tout $x$, et en
        déduire $\ker(u - \lambda) = \ker(u^* - \conj\lambda)$.
  \item Montrer que les sous-espaces propres de $u$ pour des \omterm{def:b2:reduction:eigen}{valeurs
        propres} distinctes sont orthogonaux, et que l'orthogonal d'un
        sous-espace propre est $u$-stable.
  \item Conclure par récurrence que $u$ est unitairement
        \omterm{def:b2:reduction:diag}{diagonalisable}~; et réciproquement.
\end{enumerate}
\end{exercise}

\begin{exercise}[$\star$]\label{exo:b2:hermitian:9}
Un endomorphisme est \emph{\omterm{def:b2:hermitian:adjoint}{anti-hermitien}} lorsque $u^* = -u$. Montrer
que ses \omterm{def:b2:reduction:eigen}{valeurs propres} sont imaginaires pures, que $u \mapsto \iu
u$ est une bijection des \omterm{def:b2:hermitian:adjoint}{endomorphismes hermitiens} vers les
\omterm{def:b2:hermitian:adjoint}{anti-hermitiens}, et que les endomorphismes \omterm{def:b2:hermitian:adjoint}{anti-hermitiens} sont
unitairement \omterm{def:b2:reduction:diag}{diagonalisables} (\cref{exo:b2:hermitian:8}).
\end{exercise}

\begin{exercise}[$\star\star$]\label{exo:b2:hermitian:10}
(La transformation de Fourier finie) Soit $S$ le décalage \omterm{def:b2:structures:generated}{cyclique} de
$\C^n$~: $S(x_0, x_1, \dots, x_{n-1}) = (x_{n-1}, x_0, \dots,
x_{n-2})$, et $\omega = \eu^{2\iu\pi/n}$.
\begin{enumerate}
  \item Montrer que $S$ est \omterm{def:b2:hermitian:adjoint}{unitaire}, et que les vecteurs $f_k =
        \frac{1}{\sqrt n}\bigl(1, \omega^k, \omega^{2k}, \dots,
        \omega^{(n-1)k}\bigr)$, $0 \leq k < n$, forment une base
        orthonormée de vecteurs propres~: $Sf_k = \omega^{-k}
        f_k$.
  \item En déduire que toute matrice \emph{circulante} $C =
        \sum_{j=0}^{n-1} c_jS^j$ est normale, diagonalisée par la
        même base, de \omterm{def:b2:reduction:eigen}{valeurs propres} $\widehat c(k) = \sum_j
        c_j\,\omega^{-jk}$.
\end{enumerate}
\end{exercise}

\begin{exercise}[$\star\star$]\label{exo:b2:hermitian:11}
Soit $P$ un endomorphisme idempotent ($P^2 = P$) d'un \omterm{def:b2:hermitian:def}{espace
hermitien}. Montrer que $P$ est la projection \emph{orthogonale} sur
$\operatorname{im} P$ si et seulement si $P^* = P$. Donner la matrice
de la projection orthogonale sur $\C v$ ($\norm v = 1$), et
sur un sous-espace de base orthonormée $(v_1, \dots, v_k)$.
\end{exercise}

\begin{exercise}[$\star\star\star$]\label{exo:b2:hermitian:12}
(Projecteurs spectraux par interpolation) Soit $A$ hermitienne, de
\omterm{def:b2:reduction:eigen}{valeurs propres} \emph{distinctes} $\lambda_1, \dots, \lambda_p$ et de
décomposition en sous-espaces propres $E = \bigoplus_i E_i$. On
définit les polynômes de Lagrange $L_i(X) = \prod_{j\neq i}\frac{X -
\lambda_j}{\lambda_i - \lambda_j}$. Montrer que $P_i = L_i(A)$
est la projection orthogonale sur $E_i$, que $P_iP_j = 0$ pour
$i \neq j$, $\sum_i P_i = I$, et $A = \sum_i \lambda_iP_i$
(la \emph{décomposition spectrale})~; exprimer $f(A)$ pour un
polynôme $f$ quelconque en fonction des $P_i$.
\end{exercise}

\section{Problème~: Courant--Fischer, Weyl et le calcul des valeurs
propres}

\begin{problem}\label{pb:b2:hermitian:1}
Les \omterm{def:b2:reduction:eigen}{valeurs propres} d'une matrice hermitienne ne sont pas seulement
les racines d'un polynôme~: ce sont les solutions de \emph{problèmes
d'optimisation}. Ce point de vue variationnel --- \omterm{pb:b2:hermitian:1}{quotients de
Rayleigh} et \emph{théorème du min-max de
Courant--Fischer}\index{théorème de Courant--Fischer} --- rend les
\omterm{def:b2:reduction:eigen}{valeurs propres} comparables, stables et calculables, et ce problème en
récolte les moissons classiques~: les \emph{inégalités de perturbation
de Weyl}\index{inégalité de Weyl}, l'\emph{entrelacement de Cauchy},
les inégalités de trace de Schur et de Ky Fan, la monotonie de la
racine carrée matricielle, et le \omterm{def:b2:reduction:eigen}{spectre} du laplacien discret.
Partout, $A, B, E$ sont hermitiennes sur $E = \C^n$ de \omterm{def:b2:reduction:eigen}{valeurs propres}
rangées par \omterm{def:b2:structures:generated}{ordre} décroissant $\lambda_1(A) \geq \dots \geq
\lambda_n(A)$, et $R_A(x) = \frac{\langle x, Ax\rangle}{\langle x,
x\rangle}$ pour $x \neq 0$ est le \emph{quotient de
Rayleigh}\index{quotient de Rayleigh}.

\medskip
\noindent\textbf{Partie I --- \omterm{pb:b2:hermitian:1}{Quotients de Rayleigh} et min-max.}
Fixons une base propre orthonormée $(e_1, \dots, e_n)$, $Ae_i =
\lambda_ie_i$.
\begin{enumerate}
  \item Montrer que $R_A(x)$ est réel, et que
        \[
        \lambda_n \leq R_A(x) \leq \lambda_1
        \qquad (x \neq 0),
        \]
        les deux bornes étant atteintes~: $\lambda_1 = \max R_A$,
        $\lambda_n = \min R_A$.
  \item Montrer que les points critiques de $R_A$ sont exactement les
        vecteurs propres de $A$ \emph{(développer $t \mapsto R_A(x +
        tv)$ en $t = 0$ pour $v$ quelconque, puis remplacer $v$ par
        $\iu v$)}.
  \item Soit $V_k = \operatorname{Vect}(e_1, \dots, e_k)$ et
        $W_k = \operatorname{Vect}(e_k, \dots, e_n)$. Montrer
        \[
        \min_{x \in V_k\setminus\{0\}} R_A(x) = \lambda_k
        = \max_{x \in W_k\setminus\{0\}} R_A(x) .
        \]
  \item Démontrer le \emph{théorème de Courant--Fischer}~: pour $1 \leq
        k \leq n$,
        \[
        \lambda_k = \max_{\dim V = k}\;\min_{x \in
        V\setminus\{0\}} R_A(x)
        = \min_{\dim W = n-k+1}\;\max_{x \in
        W\setminus\{0\}} R_A(x)
        \]
        \emph{(pour tout $V$ de dimension $k$~: $V \cap W_k \neq
        \{0\}$ par Grassmann, donc $\min_V R_A \leq \lambda_k$~;
        la question~3 montre que la borne est atteinte)}.
  \item (Monotonie) On écrit $A \leq B$ lorsque $B - A$ est
        positive. Déduire de la question~4~: $A \leq
        B$ implique $\lambda_k(A) \leq \lambda_k(B)$ pour tout
        $k$.
\end{enumerate}

\medskip
\noindent\textbf{Partie II --- Inégalités de Weyl.}
\begin{enumerate}[resume]
  \item Montrer que des sous-espaces $V, W \subseteq \C^n$ avec $\dim V
        + \dim W > n$ s'intersectent non trivialement, et généraliser~:
        $\dim(V_1 \cap V_2 \cap V_3) \geq \dim V_1 + \dim V_2 +
        \dim V_3 - 2n$.
  \item Démontrer l'\emph{inégalité de Weyl}~: pour $i + j - 1 \leq n$,
        \[
        \lambda_{i+j-1}(A + B) \leq \lambda_i(A) +
        \lambda_j(B)
        \]
        \emph{(intersecter les sous-espaces $W_i(A)$, $W_j(B)$ et
        $V_{i+j-1}(A+B)$ de la question~3 et compter les
        dimensions)}.
  \item Définir $\vertiii{E}_2 = \max_{\norm x = 1}\norm{Ex}$
        et montrer $\vertiii E_2 = \max_k\abs{\lambda_k(E)}$ pour
        $E$ hermitienne. En déduire le \emph{théorème de perturbation
        de Weyl}~:
        \[
        \bigl|\lambda_k(A + E) - \lambda_k(A)\bigr| \leq
        \vertiii{E}_2
        \qquad (1 \leq k \leq n) :
        \]
        chaque \omterm{def:b2:reduction:eigen}{valeur propre} est une fonction $1$-lipschitzienne de la
        matrice.
  \item (Perturbations de rang un) Soit $P$ hermitienne positive
        de rang~$1$. Montrer
        \[
        \lambda_k(A) \leq \lambda_k(A + P) \leq
        \lambda_{k-1}(A) \qquad (2 \leq k \leq n),
        \]
        ainsi que $\lambda_1(A) \leq \lambda_1(A+P)$~: les
        nouvelles \omterm{def:b2:reduction:eigen}{valeurs propres} \emph{entrelacent} les anciennes.
  \item Vérifier la question~8 numériquement~: $A = \begin{pmatrix} 2 &
        \iu\\ -\iu & 2\end{pmatrix}$
        (\cref{exo:b2:hermitian:4}~: \omterm{def:b2:reduction:eigen}{spectre} $\{3, 1\}$) et
        $E = \begin{pmatrix} 0 & 1\\ 1 & 0\end{pmatrix}$
        (\omterm{def:b2:reduction:eigen}{spectre} $\{1, -1\}$)~: calculer le \omterm{def:b2:reduction:eigen}{spectre} de $A +
        E$ et les deux membres de l'inégalité.
\end{enumerate}

\medskip
\noindent\textbf{Partie III --- Entrelacement et inégalités de
trace.}
\begin{enumerate}[resume]
  \item (Entrelacement de Cauchy) Soit $B$ la sous-matrice principale
        dominante $(n-1)\times(n-1)$ de $A$. Démontrer
        \[
        \lambda_{k+1}(A) \leq \lambda_k(B) \leq \lambda_k(A)
        \qquad (1 \leq k \leq n-1)
        \]
        \emph{(voir $\C^{n-1} \subseteq \C^n$~; sur lui, $R_B$ est
        la restriction de $R_A$~; appliquer Courant--Fischer aux
        deux niveaux)}.
  \item Itérer~: pour une sous-matrice principale $B$ de taille $n - m$,
        $\lambda_{k+m}(A) \leq \lambda_k(B) \leq
        \lambda_k(A)$.
  \item (Schur) Soient $d_1 \geq d_2 \geq \dots \geq d_n$ les
        coefficients diagonaux de $A$, triés. Démontrer, pour tout
        $k$~:
        \[
        \sum_{i=1}^{k} d_i \leq \sum_{i=1}^{k}\lambda_i(A),
        \]
        avec égalité en $k = n$ (la trace)
        \emph{(les $k$ coefficients diagonaux choisis forment une
        sous-matrice principale $k\times k$~; majorer sa trace par
        la question~12)}.
  \item (Ky Fan) Démontrer
        \[
        \sum_{i=1}^{k}\lambda_i(A)
        = \max\Bigl\{\sum_{i=1}^{k}\langle x_i, Ax_i\rangle
        : (x_1, \dots, x_k) \text{ orthonormée}\Bigr\} .
        \]
  \item Vérifier les questions~11 et~13 sur
        \[
        A = \begin{pmatrix} 2 & 1 & 0\\ 1 & 2 & 1\\
        0 & 1 & 2\end{pmatrix}
        \qquad
        \bigl(\operatorname{Sp} = \{2 - \sqrt2,\ 2,\
        2+\sqrt2\}\bigr)
        \]
        face à son bloc dominant $2\times2$ (\omterm{def:b2:reduction:eigen}{spectre} $\{1,
        3\}$) et à sa diagonale.
\end{enumerate}

\medskip
\noindent\textbf{Partie IV --- L'\omterm{def:b2:structures:generated}{ordre} de Loewner.} $A \leq B$
signifie toujours $B - A$ positive~; toutes les matrices de cette
partie sont hermitiennes.
\begin{enumerate}[resume]
  \item Montrer~: $A \leq B$ implique $a_{ii} \leq b_{ii}$ pour tout
        $i$, $\operatorname{tr} A \leq \operatorname{tr} B$,
        et $C^\dagger AC \leq C^\dagger BC$ pour \emph{toute}
        matrice complexe $C$.
  \item Montrer que l'élévation au carré n'est \emph{pas} monotone~:
        pour
        \[
        A = \begin{pmatrix} 1 & 0\\ 0 & 0\end{pmatrix},
        \qquad
        B = \begin{pmatrix} 2 & 1\\ 1 & 1\end{pmatrix},
        \]
        vérifier $0 \leq A \leq B$ mais $A^2 \not\leq B^2$.
  \item Démontrer que la racine carrée \emph{est} monotone~: $0 \leq
        A \leq B$ implique $\sqrt A \leq \sqrt B$ \emph{(soit
        $\mu$ une \omterm{def:b2:reduction:eigen}{valeur propre} de $\sqrt B - \sqrt A$ de vecteur
        propre \omterm{def:b2:hermitian:adjoint}{unitaire} $v$~; calculer $\langle v, (B - A)v\rangle =
        \mu\bigl(\langle v, \sqrt B\,v\rangle + \langle v,
        \sqrt A\,v\rangle\bigr)$ et discuter)}.
  \item Démontrer que l'inversion est \emph{antitone} sur les matrices
        définies positives~: $0 < A \leq B$ implique $B^{-1} \leq A^{-1}$
        \emph{(conjuguer par \omterm{def:b2:quadratic:def}{congruence} par $A^{-1/2}$ pour se ramener
        à $I \leq M \Rightarrow M^{-1} \leq I$, qui est scalaire dans
        une base spectrale)}.
  \item Soient $A, B$ définies positives. Montrer que les
        \omterm{def:b2:reduction:eigen}{valeurs propres} de $AB$ (non hermitienne en général~!) sont
        réelles et positives, et que
        \[
        \lambda_{\max}(AB) \leq
        \lambda_{\max}(A)\,\lambda_{\max}(B)
        \]
        \emph{(conjuguer par $\sqrt A$~: $AB \sim \sqrt
        A\,B\sqrt A$)}.
\end{enumerate}

\medskip
\noindent\textbf{Partie V --- Le laplacien discret, traité.}
Soit $T_n$ la matrice tridiagonale $n \times n$ avec $2$ sur
la diagonale et $-1$ sur les deux diagonales adjacentes.
\begin{enumerate}[resume]
  \item Avec $\theta_k = \frac{k\pi}{n+1}$, vérifier que les
        vecteurs $v_k = \bigl(\sin(j\theta_k)\bigr)_{1\leq j\leq
        n}$ vérifient $T_nv_k = (2 - 2\cos\theta_k)\,v_k$
        \emph{(identité produit-somme~; vérifier les lignes de bord
        $j = 1, n$)}. En déduire~:
        \[
        \operatorname{Sp}(T_n) = \Bigl\{4\sin^2
        \frac{k\pi}{2(n+1)} : 1 \leq k \leq n\Bigr\},
        \]
        toutes simples, toutes positives~: $T_n$ est définie positive.
  \item (Un potentiel) Pour une matrice diagonale réelle $D =
        \operatorname{diag}(d_1, \dots, d_n)$, encadrer le
        \omterm{def:b2:reduction:eigen}{spectre}~: pour tout $k$,
        \[
        \lambda_k(T_n) + \min_i d_i \;\leq\;
        \lambda_k(T_n + D) \;\leq\; \lambda_k(T_n) + \max_i
        d_i .
        \]
  \item Vérifier explicitement l'entrelacement de Cauchy entre $T_3$
        et $T_2$ (\omterm{def:b2:reduction:eigen}{spectres} $\{2 \pm \sqrt2, 2\}$ et $\{1,
        3\}$), et interpréter~: $T_2$ est $T_3$ dont on a retiré une
        extrémité du chemin.
  \item Montrer que les \omterm{def:b2:reduction:eigen}{valeurs propres} extrêmes vérifient, quand $n
        \to \infty$~:
        \[
        \lambda_{\min}(T_n) = 4\sin^2\frac{\pi}{2(n+1)}
        \sim \frac{\pi^2}{(n+1)^2},
        \qquad
        \lambda_{\max}(T_n) \to 4 ,
        \]
        de sorte que le conditionnement $\kappa_n =
        \lambda_{\max}/\lambda_{\min}$ croît comme
        $\frac{4(n+1)^2}{\pi^2}$~: discrétiser une dérivée
        seconde sur une grille de plus en plus fine est
        intrinsèquement mal conditionné.
  \item Synthèse. En une phrase chacune~: (i) pourquoi c'est la
        caractérisation variationnelle, et non le \omterm{def:b2:reduction:charpoly}{polynôme
        caractéristique}, qui rend les \omterm{def:b2:reduction:eigen}{valeurs propres} stables
        (questions~8--9)~; (ii) quelles questions n'ont utilisé que
        $\lambda_1 = \max R_A$ et lesquelles ont eu besoin du
        min-max \omterm{def:b2:metric:complete}{complet}~; (iii) ce que l'\omterm{def:b2:structures:generated}{ordre} de Loewner ajoute à
        l'histoire~; (iv) où ces outils réapparaissent (analyse
        numérique des matrices de rigidité de la question~24~;
        théorie quantique des perturbations~; et, dans le volume de
        l'année~3, le principe du min-max pour les opérateurs
        autoadjoints \omterm{def:b2:metric:compact}{compacts}).
\end{enumerate}
\end{problem}
